\chapter*{Abstract}
\addcontentsline{toc}{chapter}{Abstract}

The rapid proliferation and deceptive fluency of Large Language Models (LLMs) have created urgent security and integrity concerns across digital journalism, social platforms, and educational institutions. While machine-generated text detection and automated fact-checking methodologies have progressed rapidly for isolating, high-resource languages such as English, low-resource and morphologically complex agglutinative languages remain severely neglected. In agglutinative languages such as Kazakh, standard statistical subword tokenizers (such as Byte-Pair Encoding and SentencePiece) partition intricate inflectional and derivational affix chains into arbitrary, semantically opaque byte sequences. This pathology---formalized in this dissertation as \textit{subword token inflation}---obscures underlying root morphemes, distorts sequence probability distributions, and leads to an unacceptable false-positive accusation rate on short colloquial texts ($\le 60$ characters) of up to 12.80\%. Furthermore, existing forensic pipelines conflate generative provenance with epistemic truth, erroneously treating synthetic origin and factual veracity as interchangeable dimensions.

To resolve these foundational vulnerabilities, this dissertation presents the first unified, morphologically-grounded framework reconciling generator-invariant AI text detection, document-scale forensic robustness, and evidence-grounded factual verification for the Kazakh language. First, we construct a mathematically rigorous 83-rule Finite-State Transducer (FST) modeling the full hierarchy of Kazakh morphotactics, including nominal cases, possessive markers, verbal conjugation tiers, polarity allomorphs, and Russian-Kazakh bilingual loanword inflections. We provide formal mathematical proofs and empirical validation demonstrating that FST pre-segmentation mitigates subword token inflation by up to 36.65\% ($p < 0.0001$), compelling commercial subword vocabularies to converge toward the empirical natural Kazakh morphemic limit of 2.202 tokens per word. 

Second, we incorporate this linguistic transducer into a dual-stream gated neural architecture that dynamically fuses pretrained contextual representations (KazRoBERTa, XLM-RoBERTa) with explicit morphological feature vectors. To overcome generator-specific overfitting, we regularize the latent representation space using multi-generator Supervised Contrastive Learning (SupCon). Across our newly curated 12,848-sample KazAI-Detect benchmark and multi-domain corpora, our framework achieves a +42.18 percentage point gain in out-of-distribution AUROC, demonstrates 100.00\% AUROC on held-out Qwen-2.5-7B generations under Leave-One-Generator-Out (LOGO) cross-validation, and slashes false positive accusations on short texts by 43.21\% ($p = 0.000004$). For long-document forensic ingestion, we design a streaming \texttt{SentencePreservingChunker} fortified with 10 Kazakh abbreviation guards and dynamic Top-K worst-case pooling, maintaining an efficient memory footprint strictly below 1.4 GB VRAM when analyzing 25,000-word manuscripts.

Third, to decouple machine origin from factual veracity, we introduce \textbf{Kazakh-FEVER 3K}, the first human-validated fact verification benchmark for Kazakh comprising 3,000 claim-evidence pairs curated over a 250,000-passage encyclopedic corpus. To neutralize superficial lexical shortcuts, we formulate an adversarial \textsc{Hard NEI} (Not Enough Information) curation protocol that demands fine-grained morphological reasoning over negation polarity, temporal intervals, and evidential markers. We construct an end-to-end verification pipeline coupling a hybrid evidence retriever (morpho-stemmed BM25 combined with dense multilingual \textsc{mContriever} representations via Reciprocal Rank Fusion, achieving 99.8\% Recall@5 and 0.994 MRR) with a 16-dimensional morphological diagnostic cross-encoder. Our verifier attains 82.6\% accuracy, 82.1\% Macro-F1, and 71.4\% Strict Joint FEVER score, outperforming competitive cross-lingual baselines by +10.8 percentage points.

Finally, we synthesize detection likelihood and factual veracity into a continuous 2D Cartesian Trust Matrix $(T_{\text{fact}}, T_{\text{gen}})$, mapping document claims into four operational risk quadrants with a continuous scalar confidence metric $\mathcal{T}(x) \in [0, 1]$. The entire framework is operationalized through an open-source, interactive four-tab explainability dashboard deployed publicly on Hugging Face Spaces with sub-1.2-second cold-start latency and validated by 506 automated regression tests. This work establishes an authoritative linguistic and algorithmic foundation for defending information integrity across agglutinative language ecosystems.

\vspace{0.8em}
\noindent\textbf{Keywords:} AI-generated text detection, agglutinative morphology, finite-state transducers, Kazakh NLP, fact verification, Kazakh-FEVER, 2D trust matrix, contrastive learning, subword token inflation, out-of-distribution robustness.

\clearpage

\section*{Аңдатпа}
\addcontentsline{toc}{chapter}{Аңдатпа (Kazakh Abstract)}

Үлкен тілдік модельдердің (LLM) қарқынды дамуы мен олар генерациялаған мәтіндердің жоғары деңгейдегі табиғилығы цифрлық медиа, әлеуметтік желілер және білім беру саласында ақпараттық қауіпсіздік пен мәтіндік түпнұсқалықты сараптау тұрғысынан бұрын-соңды болмаған күрделі мәселелер туындатып отыр. Қазіргі таңда машиналық мәтіндерді анықтау және олардың дәйектілігін тексеру әдістері негізінен ағылшын тілі сияқты аналитикалық, ресурстары мол тілдерге бағытталған, ал қазақ тілі сияқты морфологиялық құрылымы бай, агглютинативті тілдердің өзіндік ерекшеліктері жеткілікті деңгейде зерттелмеген. Агглютинативті тілдерде қолданылатын стандартты статистикалық қосымша сөздіктер (BPE және SentencePiece) күрделі септік, тәуелдік және етістік жалғауларының тізбектерін мағынасыз, кездейсоқ байт бірліктеріне бөлшектеп, \textit{субсөздік инфляция} (subword token inflation) мәселесін туындатады. Бұл патология сөздердің түбір морфемаларын көмескілеп, тілдік статистиканы бұрмалайды және қысқа тұтынушылық пікірлерде ($\le 60$ таңба) жалған айыптаулар (false positives) көрсеткішін 12.80\%-ға дейін арттырады. Сонымен қатар, қолданыстағы жүйелер мәтіннің генерациялану ықтималдығы мен оның фактілік шынайылығын бір-бірімен шатастырып, синтетикалық мәтін мен жалған мәліметтің аражігін ажырата алмайды.

Аталған іргелі мәселелерді шешу үшін осы диссертациялық жұмыста қазақ тіліндегі машиналық мәтіндерді анықтау, құжаттық масштабтағы сенімділік және фактілік шынайылықты бағалаудың тұңғыш біртұтас, морфологиялық негізделген кешенді жүйесі құрылды. Біріншіден, қазақ тілінің есімдік септеулерін, тәуелдік жүйесін, етістік шақтары мен болымсыздық аффикстерін, сондай-ақ қазақ-орыс тілдік кодын ауыстыру заңдылықтарын толық қамтитын 83 ережеден тұратын Соңғы Автоматты Трансдьюсер (FST) әзірленді. Математикалық және эмпирикалық түрде FST сегментациясы субсөздік инфляцияны 36.65\%-ға төмендетіп ($p < 0.0001$), коммерциялық сөздіктерді қазақ тілінің табиғи 2.202 морфема/сөз шегіне жақындататыны дәлелденді. 

Екіншіден, осы трансдьюсер алдын ала оқытылған трансформаторлар (KazRoBERTa, XLM-RoBERTa) мен ашық морфологиялық векторларды динамикалық түрде біріктіретін қос ағынды гейттік нейрондық архитектураға енгізілді. Генераторға тәуелділікті жою үшін латентті кеңістік көп генераторлы супервайзерлік контрастивті оқытумен (SupCon) реттелді. KazAI-Detect (12,848 үлгі) бенчмаркінде жүргізілген тәжірибелер біздің моделіміздің таратылымнан тыс (OOD) деректерде AUC көрсеткішін +42.18 пайыздық тармаққа арттырып, генераторды алып тастап тексеру (LOGO) бойынша Qwen-2.5-7B моделі үшін 100.00\% нәтиже көрсеткенін және қысқа мәтіндердегі жалған айыптауларды 43.21\%-ға азайтқанын ($p = 0.000004$) растады. Ұзын мәтіндерді өңдеу үшін 10 қазақша қысқарту қорғанысымен жабдықталған \texttt{SentencePreservingChunker} және динамикалық Top-K пулдеу жүйесі әзірленіп, 25,000 сөзден тұратын монографияларды өңдеу кезіндегі бейнежад көлемі 1.4 GB VRAM деңгейінде шектелді.

Үшіншіден, мәтіннің генерациялануын оның фактілік шынайылығынан ажырату мақсатында 250,000 энциклопедиялық мәтін негізінде құрастырылған 3,000 мәлімдеме-дәлел жұбынан тұратын \textbf{Kazakh-FEVER 3K} бенчмаркі және күрделі морфологиялық алдамшы үлгілерді қамтитын \textsc{Hard NEI} хаттамасы ұсынылды. Морфо-BM25 пен тығыз \textsc{mContriever} моделін Reciprocal Rank Fusion (RRF) арқылы біріктірген гибридті іздеу жүйесі 99.8\% Recall@5 және 0.994 MRR нәтижесін берді. Ал 16 өлшемді морфологиялық кросс-энкодер 82.6\% дәлдік, 82.1\% Macro-F1 және 71.4\% қатаң бірлескен FEVER көрсеткішін қамтамасыз етті.

Соңында, генерациялық ықтималдық пен фактілік шынайылықты біріктіретін үздіксіз 2D Декарттық Сенім Матрицасы $(T_{\text{fact}}, T_{\text{gen}})$ құрылып, тұжырымдар төрт қауіптілік квадрантына жіктелді. Барлық құралдар интерактивті төрт қосымшалы Gradio жүйесі түрінде Hugging Face Spaces платформасына орналастырылып, 506 автоматты тест арқылы тексерілді.

\vspace{0.8em}
\noindent\textbf{Түйін сөздер:} Жасанды интеллект мәтіндерін анықтау, агглютинативті морфология, соңғы автоматтар (FST), қазақ тілін өңдеу (NLP), фактілерді тексеру, Kazakh-FEVER, 2D сенім матрицасы, супервайзерлік контрастивті оқыту, субсөздік инфляция, таратылымнан тыс сенімділік.

\clearpage

{\zhfont
\section*{\zhfont 摘要}
\addcontentsline{toc}{chapter}{\texorpdfstring{{\zhfont 摘要} (Chinese Abstract)}{Chinese Abstract}}

大语言模型（LLM）的快速演进与广泛普及对数字内容溯源、信息真实性审核以及学术诚信构成了严峻挑战。然而，现有的生成文本检测与事实核验技术主要聚焦于英语等高资源孤立语，在以哈萨克语为代表的形态丰富型黏着语（Agglutinative Languages）上存在严重的适应性缺陷。在黏着语中，标准统计分词算法（如 BPE 与 SentencePiece）易引发严重的子词通胀（Subword Token Inflation，扩展率达 3.8 倍），将复杂的格变化、人称屈折及派生词缀任意切分为无语义的碎片字节。这种病理破坏了核心词根边界，并在短文本场景下导致高达 12.80\% 的误报率（False Positive Rate）。同时，现有检测系统往往将机器生成概率与事实真伪混为一谈，缺乏对二者正交关系的联合解耦建模。

为解决上述基础性难题，本学位论文构建了首个面向低资源黏着语（哈萨克语）的形态学驱动多领域 AI 生成文本检测与事实核验统一框架。首先，本研究形式化设计了包含 83 条形态规则的有限状态转换器（FST），全面涵盖名词格变化、动词时态与极性屈折以及俄哈混合借词形态。数学分析与实证表明，FST 预切分将子词膨胀率降低了 36.65\%（$p < 0.0001$），促使商用分词模型收敛于哈萨克语天然的语素极限（2.202 词元/词）。在此基础上，本文提出了融合预训练语言模型（KazRoBERTa, XLM-RoBERTa）与 FST 特征的双流门控神经网络，并引入多生成器有监督对比学习（SupCon）进行表征对齐。在包含 12,848 个样本的 KazAI-Detect 基准及多生成器评估中，该架构在分布外（OOD）测试集上的 AUROC 提升了 42.18 个百分点，在针对未见生成器 Qwen-2.5-7B 的留一生成器交叉验证（LOGO）中达到 100.00\% AUROC，并将短文本误报率降低了 43.21\%（$p = 0.000004$）。此外，采用保留句子完整性的分块策略（集成 10 种哈萨克语缩写防护）与动态 Top-K 池化，确保 25,000 词长文档的显存占用低于 1.4 GB VRAM。

在事实核验方面，本研究构建了哈萨克语首个人工验证的事实核验基准数据集 \textbf{Kazakh-FEVER 3K}，基于 250,000 个百科全书知识段落涵盖 3,000 对断言-证据对，并提出针对形态陷阱的对抗性 \textsc{Hard NEI} 协议。通过将词素 BM25 与多语言稠密检索器（\textsc{mContriever}）通过倒数排名融合（RRF）相结合，检索阶段实现了 99.8\% 的 Recall@5 和 0.994 的 MRR；辅以 16 维形态学诊断交叉编码器，事实判定准确率达到 82.6\%，Macro-F1 达到 82.1\%，严格联合 FEVER 得分达到 71.4\%，显著超越跨语言基线模型。

最后，本文提出了连续二维笛卡尔信任矩阵 $(T_{\text{fact}}, T_{\text{gen}})$，将断言判定映射至四个风险象限，并依托 Hugging Face Spaces 部署了具备亚 1.2 秒冷启动能力的四标签页可解释性交互平台。全套系统通过了 506 项自动化回归测试，为低资源语言的数字信息完整性防御提供了完备且可复现的理论与工程范例。

\vspace{0.8em}
\noindent\textbf{关键词：}人工智能生成文本检测，黏着语形态学，有限状态转换器 (FST)，哈萨克语自然语言处理，事实核验，Kazakh-FEVER，2D 信任矩阵，对比学习，子词通胀，分布外鲁棒性。
}
